\section{A good signature for each prescribed lower-row parameter}

% BEGIN J86S1
\paragraph{J86S1. A joint good eigenvector over the signature coefficient field.}
Fix an original member, a location $a$, its twist $i$, its side $b$
and its index $j=j(a)$. Fix a good row $\sigma$ with
$\mathrm{isot}_j(\sigma)\ne0$. Write $L=LL$ and $L_0=LL_0$.
The norm choice identifies the rational row space, finite action,
split operators and their eigenvalues before and after localization.
Here the group element $g$ acts in rational coordinates by
$\chi(g)^{-1}\rho(g)$. To see the identifications used below,
transport the norm choice's second-step data along its equality
with the member and its pinned character frame. Its local map reads
$\mathrm{locW}(a,\sigma,v)=R_w.\mathrm{locR}(\sigma,v)$ at split
places. Its expansion identifies the native and rational form spaces
by $e_F$ and intertwines their finite-group actions. The two
descriptions of the level subgroup agree by the member's group
reading. Both level character values have the same image under the
chosen embedding $L_0\to\operatorname{Frac}(O_{0,i})$ and its
inclusion in $\operatorname{Frac}(O_i)$, so they agree by
injectivity. In particular these are the subgroup and character in
the definition of the native space of level vectors.

Choose a nonzero $y\in\mathrm{isot}_j(\sigma)$. In the complex
target model the level projector fixes $By$ and the level subgroup
acts on it by the prescribed character. The member's extra
realization has an actual scalar system $u_{\sigma,k,v}\in L_0$:
on these complex vectors the normalized spherical double-coset sum
is multiplication by the specified embedding of $u_{\sigma,k,v}$.
The exceptional set may be taken to be the intersection of the
sets $S_0(i')$ for all twists $i'$. This is finite, being contained
in $S_0(i)$, and contains the extra realization's exceptional set
and the primes over $2g_0g_1$. For a split $v$ outside the
intersection choose $i'$ with $v\notin S_0(i')$. Conjugation
stability gives the same exclusion at its conjugate, and the pinned
geometric bad primes are contained in every $S_0(i')$. The rational
Hecke reading of the extra realization therefore applies with this
$i'$: the split diagonal with $k$ uniformizers and its right-coset
representatives realizes the stated normalized operator on rational
coordinates. The finite index of the hyperspecial intersection gives
a finite set of representatives. Its choice has no effect on
level-fixed vectors.

Consider simultaneously the kernels of the rational matrices for
the good projector minus identity, the level-group actions minus
their character, and all the split operators minus
$u_{\sigma,k,v}$. The rational target model has a basis which becomes
a basis of the complex Hilbert target; its realization in
multiplicity maps is injective. Thus $By$ gives a nonzero vector
in these kernels after extension to $\mathbb C$. In finite
dimension a finite subfamily of kernel equations cuts out their
whole intersection: take a basis of the span of all their matrix
rows. Matrix rank is unchanged by extension of fields. Consequently
some $x\ne0$ solves all the equations over $L$ in the same rational
model. The equality of the good projectors uses the pinned equality
of their group-algebra coefficients; the character twist in the
action cancels the character in the rational realization. In
particular $x$ belongs to the original rational good space, including
its level-character condition.

Put $K=\operatorname{Frac}(O_i)$ and write $\beta_b:L_0\to K$
for the side's coefficient reading. The small coefficient field
$K_s$ maps into $L_0$ by the member pin. Choose a $K_s$-basis of its
form which is also an $L$-basis after extension. Choose an
$L_0$-linear functional $f:L\to L_0$ taking a nonzero coordinate
of $x$ to $1$ (extend that coordinate to an $L_0$-basis).
Applying $\beta_b f$ to each coordinate gives a nonzero vector
in the $K$-form. Every endomorphism preserving the $K_s$-form
has matrices over $K_s$; the coordinate map therefore intertwines
it with its base change. It also transports the $L_0$-valued
character and the projector coefficients. Transport back by
$e_F^{-1}$ and call the resulting vector $z$. Then
\[
 z\ne0,\qquad e_G z=z,\qquad u z=\psi_U(u)z\quad(u\in U),\qquad
 H_{k,v} z=c_{i,b,v}^{\,k}\beta_b(u_{\sigma,k,v})z .
\]
The last identity holds for $1\le k\le r$ and every split
$v\notin S_0(i)$, with the same native unit $c_{i,b,v}$ that
relates the two Hecke actions. Each $u_{\sigma,k,v}$ used here is
the scalar in the given original realization, for both sides.
% END J86S1

% BEGIN J86S2
\paragraph{J86S2. Enumeration at this row and the Bloch--Kato good vector.}
Use $z$ from J86S1. It belongs to the level-character subspace
$F_0$, and $e_Gz=z$, so it belongs to $V_0=e_GF_0$. The
native integral submodule $B_{\rm sup}$ spans the full space over
$K$. Indeed the integral-form condition in the member's realization
places a full set of generators of the rational $K_s$-form in the
span of the integral lattice. Extending scalars gives a spanning
set for the model over $K$. The native lattice comparison contains
a common nonzero integral multiple of the transported spanning
set in $B_{\rm sup}$. Dividing by this multiplier gives the
assertion over $K$. Express $z$ as a finite linear combination
of elements of $B_{\rm sup}$; clearing the denominators gives
$0\ne d\in O_i$ with $dz\in B_{\rm sup}\cap V_0$. This intersection
is exactly the image of the native lattice $B$ in the expansion.
Thus $z$ lies in the image of $K\otimes_{O_i}B$.

On the latter space the signature span in the norm choice says that
vectors belonging to the joint eigenspaces of the signatures of this
$j,b$ span. The Hecke reading of the integral embedding intertwines
all the tested operators, also after scalar extension. Hence $z$ is
a finite sum of joint eigenvectors with systems drawn from these
native signatures. If no contributing signature shared all its
eigenvalues, for each distinct system $\lambda$ in this sum choose
one test $h$ with $\lambda(h)\ne\lambda_z(h)$. The product of
$H_h-\lambda(h)$ over these choices annihilates every summand
and is nonzero scalar multiplication on $z$. Indeed each summand
is a joint eigenvector for \emph{all} factors. This contradicts
$z\ne0$. We have found a native signature $s$ with
$j_S(i,s)=j$, $\mathrm{side}(i,s)=b$ and
\[
       h_{\rm Eig}(i,s,k,v)
             =c_{i,b,v}^{\,k}\beta_b(u_{\sigma,k,v})
             \quad\hbox{in }K .
\]
With these identifications J86S1 makes $z$ itself a witness to
the definition of $\mathrm{goodS}(i,s)$: it is nonzero, fixed by
the native good projector, and has all the specified Hecke values.
Writing $\sigma_s$ for the original row read from $s$, the norm
choice's $\mathrm{rowS\_eig}$ identifies the same left hand side
with $c_{i,b,v}^{\,k}\beta_b(u_{\sigma_s,k,v})$. Cancel the unit
and use injectivity of $\beta_b$. The two rows thus have equal
coefficients at every test for this side.
% END J86S2

% BEGIN J86S3
\paragraph{J86S3. The Galois reading of the same signature and the original triple.}
Let $\Sigma$ be the list supplied in $\mathrm{HasParamW}(a,\sigma,\Sigma)$.
The native $\mathrm{rhoS}$ specification for $s$ gives
$\mathrm{isot}_{j_S(i,s)}(\sigma_s)\ne0$ and goodness of $\sigma_s$.
Since $j_S(i,s)=j(a)$, apply U86T3--4 to $\sigma_s$ at the original
$a$. After removing the length-$t$ character string in the upper
base-change list the remaining part has length one and rank $r$.
Let $\pi$ be its base twisted by $\eta$. U86T4 gives
$\mathrm{locG}(\pi,v)=\mathrm{locW}(a,\sigma_s,v)$ at
\emph{every} split $v$. Its base is conjugate self-dual by the
base-change condition, and $\eta$ is conjugate self-dual by S86A2.
It is regular algebraic: its untwisted infinity parameters form a
submultiset of the upper list, whose exponents are distinct and lie
in $(n_V-1)/2+\mathbb Z$. The length-one part adds no Speh shifts.
The infinity exponent of $\eta$ is $1/2$ or $-1/2$ at each complex
embedding. Since $n_V$ is even and $r$ odd, shifting gives distinct
exponents in $(r-1)/2+\mathbb Z$. This proves the required test at
every embedding.

At a split place outside the finite exceptional sets the local
component of each row is spherical. The scalar for the $k$-th
hyperspecial double coset is $q_v^{k(r-k)/2}$ times the $k$-th
elementary symmetric polynomial in the Satake roots; the factor is
the square root of the modulus character of the standard Levi on
that diagonal. The measure gives the hyperspecial subgroup volume
one. The member's rational normalization multiplies each root by
$\mu_{0,j}(v)^{-1}q_v^{-n_V/2}$, and the local Frobenius reading
uses inverses of the Satake roots. Equivalently, if $A$ is the
multiset of Frobenius eigenvalues in the row reading, the
embedded rational scalar is
\[
 q_v^{k(r-k)/2}e_k\bigl(\{\alpha^{-1}\mu_{0,j}(v)^{-1}
                         q_v^{-n_V/2}:\alpha\in A\}\bigr).
\]
This follows by applying the normalized double-coset operator to a
nonzero vector of that type, using the hyperspecial fixed-vector
calculation, and cancelling the vector in the Hilbert reading.
The multisets for $\sigma,\sigma_s$ have the same size $r$: their
unramified parameters come respectively from the given rank-$r$
list and the rank-$r$ cusp $\pi$. J86S2 now equates all elementary
symmetric polynomials, since each displayed norm factor is nonzero.
Equality of monic polynomials identifies their root multisets. The
map $\alpha\mapsto\alpha^{-1}\mu_{0,j}(v)^{-1}q_v^{-n_V/2}$ is
injective on these nonzero roots. The two unramified local parameters
are therefore equal. In particular we can use the cusp $\pi$ on
the signature's row and the prescribed list on the input row even if
their finite type labels differ.

Use the reading at $a$ of $\operatorname{Frac}(O_{0,i})$ in
$\overline{\mathbb Q}_\ell$. Explicitly, extend the coefficient
map $\mathrm{twist}_2(a)$ to fractions, compose with the location's
embedding into $\overline{\mathbb Q}_\ell$, and use the corresponding
isomorphism to $\mathbb C$. The integral and side-specific
$L_0$ equalities in the localized norm choice give the identities
on $\mathbb Z_\ell$ and on $L_0$ required by $\mathrm{rhoS}$;
a fraction-field identity is checked on numerators and denominators.
Apply that specification to $\pi$. It gives semisimplicity and the
Frobenius characteristic polynomial for the twist
$\pi\otimes\eta^{-1}$ with the fixed norm shift $(1-n_V)/2$.
Semisimplicity descends through this characteristic-zero field
embedding: an invariant subspace with a complement after extension
has an equivariant retraction there. The retraction equations are
linear over the original field, and a finite spanning set of the
equations suffices, so the retraction exists over that field.

At almost every split place apply the equality of local parameters
just proved and the split-place test in $\mathrm{HasParamW}$. Twisting
each constituent of $\Sigma$ by $\eta^{-1}$ twists its full
local parameter by that character. The characteristic polynomial
for $\pi\otimes\eta^{-1}$ is therefore exactly the polynomial for
the prescribed twisted list with the same norm shift. Characteristic
polynomials commute with the coefficient embedding; on $O_{0,i}$
its composite into $\mathbb C$ is the pinned composite through
$\mathrm{twist}_2(a)$. The Frobenius assertion holds for every
valuation ring and every allowed Frobenius lift in the definition,
because the $\mathrm{rhoS}$ assertion has these quantifiers and
we have removed only a finite set of places. With the index and
side equalities from J86S2 these are all clauses of
$\mathrm{IsTriple}(a,s,\Sigma)$. The same $s$ has J86S2's
$\mathrm{goodS}$ witness. This proves the original field
$\mathrm{trip}$.
% END J86S3
